\begin{reviewcomment}{3}
	The problem of attribute-origin requires further investigation.
	The issuer-signature mechanism is a placeholder for consistency, based on Poseidon, and thus does not provide cryptographic unforgeability as it stands.
	Thus, the proof proves the consistency of the given witness but does not prove that the private attributes are derived from an authorized issuer.
	This difference is particularly important in a healthcare authorization system.
	The authors should clarify which authorization properties are guaranteed by the prototype they implemented, and which are dependent on external trust.
	Even an analytical assessment of a concrete design for the integration of real credential verification would greatly improve the work.
\end{reviewcomment}

\begin{reviewresponse}{3}
	We thank the reviewer for emphasizing this distinction.
	We agree that witness consistency alone is insufficient to establish attribute origin, and that a healthcare authorization prototype must separate cryptographic credential verification from trust in the real-world issuer.

	During the revision, we clarified and documented the implemented in-circuit EdDSA-Poseidon verification path for a credential carrying fixed DUO-relevant claims.
	For the private, issuer-signed claim bits $c_0,\ldots,c_6$ for, respectively, research use, health or medical use, methods development, genetic study, ancestry, commercial use, and non-profit use, let $d_{\mathrm{DUO}}$ denote the fixed credential schema/domain tag. The circuit reconstructs
	\begin{equation*}
		m_{\mathrm{cred}}=\mathsf{Poseidon}(d_{\mathrm{DUO}}\parallel h\parallel
		\mathit{id}_{\mathrm{req}}\parallel\mathit{id}_{\mathrm{res}}\parallel
		c_0\parallel\cdots\parallel c_6\parallel
		t_{\mathrm{from}}\parallel t_{\mathrm{until}}\parallel\gamma)
	\end{equation*}
	Before proof generation, an issuer computes $\sigma\leftarrow\mathsf{Sign}_{\mathrm{EdDSA\text{-}Poseidon}}(s_{\mathrm{iss}},m_{\mathrm{cred}})$ with its private signing key $s_{\mathrm{iss}}$.
	The resulting credential is supplied to the TPS as private off-chain evidence for proof generation.
	At proof generation, the TPS supplies the credential as private witness and the circuit verifies the issuer's EdDSA-Poseidon signature on this message.
	It also constrains the supplied issuer point to the issuer-key commitment contained in the policy-bound parameters, enforces the credential validity interval and a non-zero credential identifier, and, for a policy that activates this holder-bound credential slot, binds the requester pseudonym to a holder secret.
	The circuit derives the DUO terms $\mathrm{GRU}$, $\mathrm{HMB}$, $\mathrm{NMDS}$, $\mathrm{GSO}$, $\mathrm{NPUNCU}$, $\mathrm{NPOA}$, $\mathrm{NPU}$, and $\mathrm{NCU}$ from these signed claims.
	Under the unforgeability of EdDSA for correctly generated Baby Jubjub keys and the collision resistance of Poseidon, a TPS cannot modify the signed policy, requester, resource, claims, validity interval, or credential identifier and still produce an accepting proof.

	We have also narrowed the claim carefully.
	The prototype proves signature validity under a policy-selected issuer key; it does not prove that this key belongs to an accredited hospital, ethics board, or other healthcare authority.
	The provider currently selects the issuer-key commitment, and there is no independently governed issuer registry or credential-revocation check.
	Ontology predicates prove that a supplied code belongs to a provider-committed set, while requester allowlists prove that a holder pseudonym belongs to a provider-committed list; neither mechanism independently authenticates every associated real-world fact.
	Furthermore, the current TPS receives the holder secret and credential witness, so holder binding does not protect against a malicious TPS reusing otherwise valid evidence within its signed scope.
	The Polygon Amoy evaluation uses a deterministic test-only issuer key to exercise the cryptographic path and measure the complete system; it is not evidence of institutional issuance or production key management.

	We added an analytical production-integration design to the Discussion.
	A governed issuer registry would publish schema-qualified, time-bounded issuer keys and require canonical encoding, subgroup validation, and proof of possession at registration.
	The issuer would sign a domain-separated credential containing a holder-key commitment, claims, authorization scope, validity interval, and credential identifier.
	The circuit would verify the credential signature and private membership of the issuer key in the active registry.
	The holder would retain its secret and sign the authorization context for each session, preventing the TPS from creating fresh-session proofs from a captured credential presentation.
	An epoch-specific sparse-Merkle or accumulator root would additionally support private non-revocation proofs.
	These roots can be folded into a fixed-format state commitment, preserving the policy-independent public interface.
	For credentials using Ed25519, ECDSA, or another non-circuit-friendly signature, the alternatives are to pay the higher in-circuit verification cost or have the issuer additionally issue a SNARK-friendly credential; TPS re-signing would not preserve issuer-origin assurance.

	The implemented credential path is covered by the DUO profile and tamper suites.
	The targeted tests accept a correctly signed credential and reject a non-committed issuer, a forged signature under the committed key, credentials outside their validity interval, substitutions of the policy, requester, resource, or signed claims, inconsistent derived attributes, and an invalid holder binding.
	All 25 targeted tests pass, including full Groth16 proof and on-chain grant/deny cases, and the frozen manifest, verifier ABI, and cryptographic artifacts validate successfully.

	\textbf{Changes in the manuscript.}
	The Problem and Threat Model now separates issuer-signed profile claims from ontology and allowlist evidence and states the remaining issuer-governance, revocation, identity, and TPS trust assumptions.
	The Design section replaces the obsolete placeholder description with the implemented credential relation and its exact scope.
	The Discussion distinguishes cryptographic signature validity from real-world issuer authorization and presents a concrete registry, holder-proof, and revocation integration design.

	\RevisionLocation{Section ``Problem and Threat Model,'' subsection ``System and Trust Model.''}
	\begin{revisiontext}
		Attribute-origin guarantees depend on the evidence type.
		For profile-backed DUO terms, the prototype verifies an EdDSA-Poseidon credential under a provider-selected, policy-bound issuer-key commitment.
		The credential binds the policy commitment, requester pseudonym, protected resource, claims, validity interval, and credential identifier; the corresponding policy attributes are derived inside the circuit.
		Thus, subject to signature unforgeability and Poseidon collision resistance, these values cannot be modified without invalidating the proof.
		This guarantee does not establish the real-world authority of the issuer or the truthfulness of issuance.
		Ontology and requester-allowlist predicates likewise establish membership only relative to provider-committed sets.
		Issuer-key governance, credential revocation, real-world identity verification, and protection against a compromised TPS remain external trust and lifecycle functions.
	\end{revisiontext}

	\RevisionLocation{Section ``TPS-ZKPI Design,'' subsection ``Credential-Origin Binding.''}
	\begin{revisiontext}
		Let $c_0,\ldots,c_6$ be the private, issuer-signed claim bits for research use, health or medical use, methods development, genetic study, ancestry, commercial use, and non-profit use, respectively.
		Let $d_{\mathrm{DUO}}$ be the fixed credential schema/domain tag.
		The circuit derives the DUO terms $\mathrm{GRU}$, $\mathrm{HMB}$, $\mathrm{NMDS}$, $\mathrm{GSO}$, $\mathrm{NPUNCU}$, $\mathrm{NPOA}$, $\mathrm{NPU}$, and $\mathrm{NCU}$ from these claims.
		For issuer public key $A=(A_x,A_y)$, the circuit reconstructs the domain-separated message
		\begin{equation*}
			m_{\mathrm{cred}}=\mathsf{Poseidon}(d_{\mathrm{DUO}}\parallel h\parallel
			\mathit{id}_{\mathrm{req}}\parallel\mathit{id}_{\mathrm{res}}\parallel
			c_0\parallel\cdots\parallel c_6\parallel
			t_{\mathrm{from}}\parallel t_{\mathrm{until}}\parallel\gamma)
		\end{equation*}
		Before proof generation, an issuer holding the corresponding private signing key $s_{\mathrm{iss}}$ issues the credential by computing $\sigma\leftarrow\mathsf{Sign}_{\mathrm{EdDSA\text{-}Poseidon}}(s_{\mathrm{iss}},m_{\mathrm{cred}})$.
		The resulting credential is supplied to the TPS as private off-chain evidence for proof generation.
		At proof generation, the TPS supplies $A$, $\sigma$, and the credential fields as private witness values.
		The circuit then enforces
		\begin{equation*}
			\mathsf{Verify}_{\mathrm{EdDSA\text{-}Poseidon}}(A,\sigma,m_{\mathrm{cred}})=1,
			\qquad
			\mathsf{Poseidon}(A_x\parallel A_y)=k_{\mathrm{iss}}.
		\end{equation*}
		The circuit additionally requires
		\begin{equation*}
			t_{\mathrm{from}}\leq\tau\leq t_{\mathrm{until}},
			\qquad \gamma\neq0,
		\end{equation*}
		For a policy that activates this holder-bound credential slot, it also requires $\mathit{id}_{\mathrm{req}}=\mathsf{Poseidon}(s_{\mathrm{holder}})$, binding the credential to the pseudonymous requester identifier in the authorization context.
		The corresponding DUO attributes are constrained to these claim-derived values rather than accepted independently from the prover.
		This is a cryptographic signature check, not a hash-consistency placeholder.
	\end{revisiontext}

	\RevisionLocation{Section ``Privacy Evaluation,'' subsection ``Dataset and corpus generation.''}
	\begin{revisiontext}
		For reproducibility, the evaluation generator uses a deterministic test issuer to create the private profile-credential witnesses used by policies that exercise the credential predicate; this testing key does not model issuer governance or production key management.
	\end{revisiontext}

	\RevisionLocation{Section ``Discussion.''}
	\begin{revisiontext}
		The remaining trust anchor is therefore issuer governance.
		The prototype binds the proof to a policy-selected issuer-key commitment, but it has no independent registry establishing that this key belongs to an accredited healthcare authority, and it does not check credential revocation.
		A concrete deployment can close this gap without changing the Boolean interpreter.
		A governed issuer registry would publish schema-qualified issuer keys and their validity epochs; the circuit would verify both the credential signature and a Merkle proof that the issuer key is active in the registry.
		The holder would retain its secret key and authorize each request by signing the context digest and request epoch, while an epoch-specific sparse-Merkle or accumulator root would support a private non-revocation witness for the credential identifier.
	\end{revisiontext}
\end{reviewresponse}
